\chapter{The Argument in Summary}
\label{c:summary}

\begin{glance}
The question is not who builds the better model: the measured gap at the top is a few index points---eight on the re-based index~\cite{aa-v43}---and has been narrow for over a year. The question is where the \emph{rent} lands when intelligence becomes cheap, because all three precedent campaigns were won that way and none with a better artifact. The answer comes in three parts. The strategy is real, coherent and further along than a layer-by-layer reading suggests. The method is no longer Chinese. And the outcome is not settled, because the layers that stay scarce are still mostly Western-held, and a West that commoditizes competently captures the surplus of a cheap-intelligence world as surely as a China that does. \full{29}
\end{glance}

\section{The answer in three parts}

\textbf{First, the strategy is real and further along than a layer-by-layer reading suggests.} Every one of the playbook's nine moves is occupied in the AI campaign in a recognizable form and at a datable stage (Chapter~\ref{c:assembled}): not a plan on paper but an execution in progress.

\textbf{Second---and this is the finding that most changed the book---the method is no longer Chinese.} Over 2026 the largest Western firms converged on the same instrument, each commoditizing a layer it does not own to defend the layer it does (Chapter~\ref{c:mirror}): the accelerator vendor gives away weights with their training data and has now bought the platform and the people that make them. That strengthens Chapter~\ref{c:precedents}---a playbook that predicts the conduct of actors it was not derived from is doing more work than one that does not---and complicates every score in Chapter~\ref{c:elements}, because a layer being commoditized by two parties at once tells you less about either of them.

\textbf{Third, the outcome is not settled, and this book's own thesis is the claim most at risk from its own findings.} If every layer is cheap the surplus accrues to whoever holds the layer that stays scarce, and at \datafreeze\ those---the installed programming model, the enterprise deployment surface, trust and certification, deployment-generated data---are still mostly Western-held. The position is therefore conditional: China is executing the strategy that won three technology wars, it is executing it well, and whether that wins a fourth depends on where the rent lands, which is not yet observable.

\begin{table}[htbp]
\centering\scriptsize
\caption{The four propositions of Section~\ref{c:propositions} at \datafreeze: the grade the evidence supports, not the probability the author assigns \fullapp{J}.}
\label{tab:summary-props}
\begin{tabular}{@{}p{7mm}p{38mm}p{63mm}p{31mm}@{}}
\toprule
& \textbf{Claim} & \textbf{Where the evidence stands} & \textbf{Status} \\
\midrule
P1 & Open weights become the default for deployed intelligence & Majority of routed tokens; a ratchet that survived a Western re-entry; the gap at the top eight index points on the re-based board; volume on the open sparse tier, spend in the closed frontier & Distribution done; monetization contested \\
P2 & China builds a sovereign, capacity-first stack to serve it & Open ISA qualified as a CUDA host; a GPU-free machine at No.~1 on HPL and HPCG; production serving on a domestic cluster; HBM still binding, the training test not passed & Trending; integration test outstanding \\
P3 & The rent financing the Western frontier is destroyed or relocated & Mechanism overdetermined---the incumbent supplier funds open weights itself---but its datacentre revenue grew 117\% in this quarter, and its cheapest tier undercut the commodity and took share & Mechanism operating; magnitude contested \\
P4 & The position converts into durable strategic advantage & Endgame begun on consolidation, not on gatekeeping; the scarce layers still mostly Western-held; deployment-generated data and trust unresolved & Pending; least evidenced \\
\bottomrule
\end{tabular}
\end{table}

The coupling matters more than any row. P1 without P2 is adoption with no hardware base to serve it, which the precedents say decays; P2 without P1 is a sovereign stack with nothing running on it; P3 without P4 is a price war that destroys a rent pool and hands it to nobody---live now that both sides are destroying each other's rents at once. \emph{P4 is where the title's claim rests and where the argument is weakest}; attack it there rather than at the benchmarks.

\section{What would change the answer}

Six observations would move the author's position furthest before the end of 2027; each is in the dashboard of Appendix~\ref{c:register}.

\begin{enumerate}\setlength{\itemsep}{0pt}\setlength{\parskip}{0pt}
\item \textbf{A frontier-class model trained end to end on domestic silicon, packaging and memory}---the one integration test that moves P2 from trending to done. Production-scale \emph{serving} has been claimed, and a leading laboratory has ordered a year of domestic accelerators for inference while keeping training on NVIDIA; the one frontier-scale domestic-training claim in circulation has no locatable primary, and that distinction is the whole of the remaining gap.
\item \textbf{Merchant-GPU datacentre revenue growth decelerating while merchant custom-ASIC revenue holds its guided trajectory}, which confirms the silicon-commoditization pincer of Chapter~\ref{c:mirror}. The contrary observation---custom accelerators converging onto the incumbent's interconnect and memory, as one already has---confirms absorption instead, which currently has the better of the evidence.
\item \textbf{A sustained fall in closed-frontier gross margin, disclosed}: P3's transmission on the line that matters. One cut to a cheapest tier, with the flagship held and the counterparty raising prices in the same fortnight, is not that.
\item \textbf{Enterprise open-weight \emph{spend} share turning upward.} The spend share fell during the period capability reached parity; if it turns, the deployment-layer objection weakens sharply.
\item \textbf{Gatekeeping exercised on any AI layer}---the endgame move with a template in batteries and rare earths, not yet used here, whose use would settle P4 quickly and in one direction.
\item \textbf{An audited instance of a research loop closing}---a model materially accelerating the science that produces its successor, or the silicon that runs it. This would make the recursive-improvement counter-thesis of Chapter~\ref{c:convergence} decisive against everything else in this book, and it is the falsifier the author takes most seriously.
\end{enumerate}

\section{What this book is not claiming}

Four disclaimers, because a book with this title attracts readings it does not support. It is \textbf{not claiming that China has won}, or that victory is inevitable, but that a recognizable strategy is being executed competently against a stack whose rent layers are being dismantled from both sides, and that the outcome turns on an ownership question not yet observable.

It is \textbf{not claiming that the West is passive}. The evidence of 2026 is the opposite: the incumbent priced underneath the commodity at the bottom of its ladder and took share doing it~\cite{price-inversion}, grew datacentre revenue 117\% in the quarter this edition closes in~\cite{nvda-q2fy27}, and runs the same commoditization strategy with more capital behind it than any Chinese actor commands.

It is \textbf{not claiming that openness is a Chinese property} or a Chinese virtue. Openness is an instrument, available to anyone with an adjacent layer to defend, and the largest single act of open-model funding in the period this book covers was American~\cite{nvda-poolside}.

And it is \textbf{not claiming more confidence than the evidence carries}. The load-bearing weaknesses, without softening: several mid-2026 details rest on secondary aggregators; the capacity-first architecture of Chapter~\ref{c:elements} is a modeled design with no measured realization; the compute sizing is measurement-anchored for one country only; and the judgment that breadth compounds faster than depth is the claim most likely to be shown wrong. Both editions are living documents, and the register of Appendix~\ref{c:register} exists so that this argument can be scored rather than believed.
